\documentclass[12pt]{article}
\usepackage[utf8]{inputenc}
\usepackage[portuguese]{babel}
\usepackage{geometry}
\geometry{a4paper, margin=2cm}
\usepackage{courier}
\usepackage{float}
\usepackage{graphicx}
\usepackage{tabularx}
\usepackage{booktabs}
\usepackage[most]{tcolorbox}
\tcbuselibrary{skins}

% Ambiente customizado para quadros de comandos
\newtcolorbox{cmdbox}[1][]{
  colback=blue!5!white,
  colframe=blue!75!black,
  fonttitle=\bfseries,
  title=#1,
  arc=4pt,
  boxrule=0.5pt,
  enhanced,
  breakable,
  fontupper=\small,
  before upper={\setlength{\extrarowheight}{4pt}}
}

% Coluna de comando com fonte monoespaçada
\newcolumntype{C}{>{\ttfamily}l}

% Espaçamento extra nas linhas dos quadros
\newcommand{\cmdboxspacing}{\setlength{\extrarowheight}{4pt}}

\begin{document}

\begin{center}
\textbf{\Large Material de apoio ao módulo 2 - metrologia por imagem: automação com systemd}

\vspace{0.5cm}

\textit{Prof. Felipe Walter Dafico Pfrimer}

\vspace{1cm}
\end{center}

\section{Objetivo do Módulo}

No módulo anterior, vimos como acessar um sistema Linux por SSH, navegar pelo terminal, manipular arquivos e criar scripts Bash simples. Neste módulo, vamos usar esses conhecimentos para automatizar a execução de tarefas usando o \texttt{systemd}.

O objetivo principal é compreender como transformar uma rotina manual de captura de imagens em um processo controlado pelo sistema operacional. Para isso, vamos estudar três elementos:

\begin{cmdbox}[Elementos da Automação]
\begin{tabularx}{\linewidth}{@{}C >{\raggedright\arraybackslash}X >{\raggedright\arraybackslash}X@{}}
\toprule
\textbf{Elemento} & \textbf{Função} & \textbf{Exemplo no Módulo} \\
\midrule
\texttt{script} & Define o que será feito & Capturar uma imagem com \texttt{fswebcam} \\
\texttt{service} & Define como a tarefa será executada pelo sistema & Executar \texttt{captura.sh} \\
\texttt{timer} & Define quando a tarefa será executada & Repetir a captura a cada intervalo \\
\bottomrule
\end{tabularx}
\end{cmdbox}

Ao final da aula, o aluno deverá ser capaz de criar um script de captura, associá-lo a um serviço do \texttt{systemd}, repetir sua execução com um temporizador e consultar o estado e os logs da automação.

\section{Por que Automatizar a Captura de Imagens?}

Em um experimento de metrologia por imagem, a forma como os dados são coletados influencia diretamente a confiabilidade dos resultados. Se uma pessoa precisa lembrar de executar um comando a cada intervalo de tempo, podem ocorrer atrasos, esquecimentos ou variações no procedimento.

A automação reduz essa dependência da intervenção manual. Quando o sistema operacional executa a tarefa automaticamente, a rotina tende a ser mais regular, mais fácil de verificar e mais adequada para experimentos longos.

\begin{cmdbox}[Benefícios da Automação]
\begin{tabularx}{\linewidth}{@{}C >{\raggedright\arraybackslash}X >{\raggedright\arraybackslash}X@{}}
\toprule
\textbf{Aspecto} & \textbf{Problema Manual} & \textbf{Com Automação} \\
\midrule
Tempo & O operador pode atrasar ou esquecer & O sistema executa no intervalo definido \\
Repetibilidade & Cada execução pode variar & A mesma rotina é repetida \\
Registro & Erros podem passar despercebidos & O serviço pode gerar logs \\
Escala & Coletas longas cansam o operador & O sistema pode operar por horas ou dias \\
\bottomrule
\end{tabularx}
\end{cmdbox}

Automatizar não significa apenas ``fazer o computador repetir uma tarefa''. Significa transformar uma rotina experimental em um procedimento mais controlado, verificável e reprodutível.

\section{O que é systemd}

O \texttt{systemd} é um sistema de inicialização e gerenciamento de serviços usado por muitas distribuições Linux. Ele é responsável por iniciar, controlar e monitorar processos importantes do sistema, como rede, login, SSH, sincronização de horário e rotinas definidas pelo administrador.

Na prática, o \texttt{systemd} permite que uma tarefa seja gerenciada pelo próprio sistema operacional, em vez de depender de um terminal aberto pelo usuário.

\subsection{Unidades do systemd}

O \texttt{systemd} organiza suas configurações em arquivos chamados unidades. Cada unidade descreve uma tarefa, um recurso ou um grupo de ações.

\begin{cmdbox}[Tipos Comuns de Unidade]
\begin{tabularx}{\linewidth}{@{}C >{\raggedright\arraybackslash}X >{\raggedright\arraybackslash}X@{}}
\toprule
\textbf{Tipo} & \textbf{Função} & \textbf{Exemplo} \\
\midrule
\texttt{.service} & Define um serviço ou tarefa executável & \texttt{captura.service} \\
\texttt{.timer} & Agenda a execução de um serviço & \texttt{captura.timer} \\
\texttt{.target} & Agrupa estados ou etapas do sistema & \texttt{multi-user.target} \\
\bottomrule
\end{tabularx}
\end{cmdbox}

Neste módulo, vamos trabalhar principalmente com arquivos \texttt{.service} e \texttt{.timer}.

\section{O comando systemctl}

O comando \texttt{systemctl} é a principal forma de controlar unidades do \texttt{systemd}. Com ele, podemos iniciar serviços, parar serviços, verificar estados, habilitar execução automática e listar temporizadores.

\begin{cmdbox}[Comandos Básicos com \texttt{systemctl}]
\begin{tabularx}{\linewidth}{@{}C >{\raggedright\arraybackslash}X >{\raggedright\arraybackslash}X@{}}
\toprule
\textbf{Comando} & \textbf{Função} & \textbf{Exemplo de Uso} \\
\midrule
\texttt{start} & Iniciar uma unidade & \texttt{sudo systemctl start captura.service} \\
\texttt{stop} & Parar uma unidade & \texttt{sudo systemctl stop captura.timer} \\
\texttt{status} & Verificar o estado & \texttt{systemctl status captura.service} \\
\texttt{enable} & Habilitar inicialização automática & \texttt{sudo systemctl enable captura.timer} \\
\texttt{disable} & Desabilitar inicialização automática & \texttt{sudo systemctl disable captura.timer} \\
\texttt{daemon-reload} & Recarregar arquivos modificados & \texttt{sudo systemctl daemon-reload} \\
\texttt{list-timers} & Listar temporizadores ativos & \texttt{systemctl list-timers} \\
\bottomrule
\end{tabularx}
\end{cmdbox}

\textbf{Observação}: sempre que um arquivo \texttt{.service} ou \texttt{.timer} for criado ou modificado, é necessário executar \texttt{sudo systemctl daemon-reload}. Esse comando informa ao \texttt{systemd} que ele deve recarregar suas configurações.

\section{Teste Manual Antes da Automação}

Antes de criar um serviço ou um timer, é essencial testar a captura manualmente. Se a câmera, o comando ou o script não funcionam no terminal, também não funcionarão corretamente quando forem chamados pelo \texttt{systemd}.

\begin{cmdbox}[Verificação Manual da Câmera]
\begin{tabularx}{\linewidth}{@{}C >{\raggedright\arraybackslash}X >{\raggedright\arraybackslash}X@{}}
\toprule
\textbf{Comando} & \textbf{Função} & \textbf{Exemplo de Uso} \\
\midrule
\texttt{ls /dev/video*} & Verificar dispositivos de vídeo & \texttt{ls /dev/video*} \\
\texttt{type fswebcam} & Verificar se \texttt{fswebcam} existe & \texttt{type fswebcam} \\
\texttt{mkdir -p} & Criar pasta para imagens & \texttt{mkdir -p imagens} \\
\texttt{fswebcam} & Capturar imagem da webcam & \texttt{fswebcam -r 1280x720 imagens/teste.jpg} \\
\texttt{ls -lh} & Conferir o arquivo gerado & \texttt{ls -lh imagens} \\
\bottomrule
\end{tabularx}
\end{cmdbox}

Uma regra prática importante é:

\begin{center}
\textbf{Se não funciona manualmente, ainda não está pronto para virar automação.}
\end{center}

\section{Script de Captura Simples}

O primeiro modelo de automação usa um script que faz apenas uma captura e termina. Esse formato é adequado para ser chamado por um timer, pois cada disparo executa uma nova coleta.

\subsection{Exemplo de Script}

O arquivo pode ser chamado de \texttt{captura.sh}. O exemplo abaixo salva a imagem em uma pasta chamada \texttt{imagens}, usando data e hora no nome do arquivo.

\begin{cmdbox}[Script \texttt{captura.sh}]
\begin{tabularx}{\linewidth}{@{}C >{\raggedright\arraybackslash}X@{}}
\toprule
\textbf{Linha} & \textbf{Código} \\
\midrule
1 & \texttt{\#!/bin/bash} \\
2 & \texttt{} \\
3 & \texttt{mkdir -p /home/g0/imagens} \\
4 & \texttt{} \\
5 & \texttt{ts=\$(date +\%Y\%m\%d\_\%H\%M\%S)} \\
6 & \texttt{arquivo="/home/g0/imagens/foto\_\$\{ts\}.jpg"} \\
7 & \texttt{} \\
8 & \texttt{echo "Capturando imagem em \$arquivo"} \\
9 & \texttt{/usr/bin/fswebcam -r 1280x720 "\$arquivo"} \\
10 & \texttt{echo "Captura finalizada"} \\
\bottomrule
\end{tabularx}
\end{cmdbox}

O usuário \texttt{g0} deve ser substituído pelo usuário correto de cada grupo, como \texttt{g1}, \texttt{g2}, \texttt{g3} ou \texttt{g4}. Em arquivos usados pelo \texttt{systemd}, é recomendável usar caminhos absolutos, como \texttt{/home/g0/captura.sh}, em vez de caminhos relativos, como \texttt{./captura.sh}.

\subsection{Executando o Script Manualmente}

\begin{cmdbox}[Teste Manual do Script]
\begin{tabularx}{\linewidth}{@{}C >{\raggedright\arraybackslash}X >{\raggedright\arraybackslash}X@{}}
\toprule
\textbf{Comando} & \textbf{Função} & \textbf{Exemplo de Uso} \\
\midrule
\texttt{chmod +x} & Tornar o script executável & \texttt{chmod +x captura.sh} \\
\texttt{./captura.sh} & Executar a captura manualmente & \texttt{./captura.sh} \\
\texttt{ls -lh} & Verificar imagens salvas & \texttt{ls -lh imagens} \\
\bottomrule
\end{tabularx}
\end{cmdbox}

Esse teste confirma se o script consegue criar a pasta, gerar o nome do arquivo e chamar a câmera corretamente.

\section{Serviço para Captura Simples}

Um serviço do \texttt{systemd} descreve como uma tarefa deve ser executada. Neste exemplo, usaremos um serviço do tipo \texttt{oneshot}, pois a tarefa executa uma vez e termina.

\subsection{Estrutura Geral de um Serviço}

\begin{cmdbox}[Seções de um Arquivo \texttt{.service}]
\begin{tabularx}{\linewidth}{@{}C >{\raggedright\arraybackslash}X >{\raggedright\arraybackslash}X@{}}
\toprule
\textbf{Seção} & \textbf{Função} & \textbf{Exemplo} \\
\midrule
\texttt{[Unit]} & Descreve a unidade e dependências & \texttt{Description=Captura de imagem} \\
\texttt{[Service]} & Define como executar a tarefa & \texttt{ExecStart=...} \\
\texttt{[Install]} & Define como habilitar a unidade & \texttt{WantedBy=multi-user.target} \\
\bottomrule
\end{tabularx}
\end{cmdbox}

\subsection{Exemplo de \texttt{captura.service}}

\begin{cmdbox}[Arquivo \texttt{/etc/systemd/system/captura.service}]
\begin{tabularx}{\linewidth}{@{}C >{\raggedright\arraybackslash}X@{}}
\toprule
\textbf{Linha} & \textbf{Código} \\
\midrule
1 & \texttt{[Unit]} \\
2 & \texttt{Description=Captura de imagem com script do usuario g0} \\
3 & \texttt{After=network.target} \\
4 & \texttt{} \\
5 & \texttt{[Service]} \\
6 & \texttt{Type=oneshot} \\
7 & \texttt{User=g0} \\
8 & \texttt{WorkingDirectory=/home/g0} \\
9 & \texttt{ExecStart=/home/g0/captura.sh} \\
10 & \texttt{} \\
11 & \texttt{[Install]} \\
12 & \texttt{WantedBy=multi-user.target} \\
\bottomrule
\end{tabularx}
\end{cmdbox}

\begin{cmdbox}[Campos Importantes do Serviço]
\begin{tabularx}{\linewidth}{@{}C >{\raggedright\arraybackslash}X >{\raggedright\arraybackslash}X@{}}
\toprule
\textbf{Campo} & \textbf{Função} & \textbf{Neste Exemplo} \\
\midrule
\texttt{Type=oneshot} & Executa uma tarefa e termina & Uma captura por execução \\
\texttt{User=g0} & Define o usuário da execução & Usa o diretório do grupo \texttt{g0} \\
\texttt{WorkingDirectory} & Define o diretório de trabalho & \texttt{/home/g0} \\
\texttt{ExecStart} & Define o comando executado & \texttt{/home/g0/captura.sh} \\
\bottomrule
\end{tabularx}
\end{cmdbox}

\subsection{Testando o Serviço}

\begin{cmdbox}[Controle do Serviço]
\begin{tabularx}{\linewidth}{@{}C >{\raggedright\arraybackslash}X >{\raggedright\arraybackslash}X@{}}
\toprule
\textbf{Comando} & \textbf{Função} & \textbf{Exemplo de Uso} \\
\midrule
\texttt{daemon-reload} & Recarregar unidades modificadas & \texttt{sudo systemctl daemon-reload} \\
\texttt{start} & Executar o serviço & \texttt{sudo systemctl start captura.service} \\
\texttt{status} & Verificar o estado & \texttt{systemctl status captura.service} \\
\texttt{journalctl} & Ver mensagens registradas & \texttt{journalctl -u captura.service -n 20} \\
\bottomrule
\end{tabularx}
\end{cmdbox}

Depois de iniciar o serviço, a pasta de imagens deve conter uma nova captura. Se isso não acontecer, o primeiro passo é verificar o estado com \texttt{systemctl status} e depois consultar os logs com \texttt{journalctl}.

\section{Timer para Repetir a Captura}

Um timer agenda a execução de um serviço. Ele não executa a captura diretamente; ele apenas dispara o serviço no intervalo definido.

\begin{center}
\textbf{O serviço define o que fazer. O timer define quando fazer.}
\end{center}

\subsection{Estrutura Geral de um Timer}

\begin{cmdbox}[Seções de um Arquivo \texttt{.timer}]
\begin{tabularx}{\linewidth}{@{}C >{\raggedright\arraybackslash}X >{\raggedright\arraybackslash}X@{}}
\toprule
\textbf{Seção} & \textbf{Função} & \textbf{Exemplo} \\
\midrule
\texttt{[Unit]} & Descreve o temporizador & \texttt{Description=Timer de captura} \\
\texttt{[Timer]} & Define o agendamento & \texttt{OnUnitActiveSec=1min} \\
\texttt{[Install]} & Define como habilitar o timer & \texttt{WantedBy=timers.target} \\
\bottomrule
\end{tabularx}
\end{cmdbox}

\subsection{Exemplo de \texttt{captura.timer}}

\begin{cmdbox}[Arquivo \texttt{/etc/systemd/system/captura.timer}]
\begin{tabularx}{\linewidth}{@{}C >{\raggedright\arraybackslash}X@{}}
\toprule
\textbf{Linha} & \textbf{Código} \\
\midrule
1 & \texttt{[Unit]} \\
2 & \texttt{Description=Timer para executar captura de imagem periodicamente} \\
3 & \texttt{} \\
4 & \texttt{[Timer]} \\
5 & \texttt{OnBootSec=30s} \\
6 & \texttt{OnUnitActiveSec=1min} \\
7 & \texttt{Unit=captura.service} \\
8 & \texttt{} \\
9 & \texttt{[Install]} \\
10 & \texttt{WantedBy=timers.target} \\
\bottomrule
\end{tabularx}
\end{cmdbox}

Neste exemplo, o timer aguarda 30 segundos após a inicialização e, depois, executa o serviço novamente a cada 1 minuto. O intervalo de 1 minuto é útil em aula porque permite observar as imagens sendo geradas. Em uma aplicação real, o intervalo pode ser maior.

\subsection{Controlando o Timer}

\begin{cmdbox}[Comandos para o Timer]
\begin{tabularx}{\linewidth}{@{}C >{\raggedright\arraybackslash}X >{\raggedright\arraybackslash}X@{}}
\toprule
\textbf{Comando} & \textbf{Função} & \textbf{Exemplo de Uso} \\
\midrule
\texttt{enable --now} & Habilitar e iniciar imediatamente & \texttt{sudo systemctl enable --now captura.timer} \\
\texttt{status} & Verificar estado do timer & \texttt{systemctl status captura.timer} \\
\texttt{list-timers} & Listar temporizadores ativos & \texttt{systemctl list-timers} \\
\texttt{stop} & Parar temporariamente & \texttt{sudo systemctl stop captura.timer} \\
\texttt{disable} & Desabilitar inicialização automática & \texttt{sudo systemctl disable captura.timer} \\
\bottomrule
\end{tabularx}
\end{cmdbox}

Quando o timer está ativo, o serviço será chamado automaticamente. Para verificar se as capturas ocorreram, consulte a pasta de imagens e os logs do serviço:

\begin{cmdbox}[Verificação da Captura Agendada]
\begin{tabularx}{\linewidth}{@{}C >{\raggedright\arraybackslash}X >{\raggedright\arraybackslash}X@{}}
\toprule
\textbf{Comando} & \textbf{Função} & \textbf{Exemplo de Uso} \\
\midrule
\texttt{ls -lh} & Verificar arquivos gerados & \texttt{ls -lh imagens} \\
\texttt{journalctl} & Ver logs da execução & \texttt{journalctl -u captura.service -n 30} \\
\bottomrule
\end{tabularx}
\end{cmdbox}

\section{Logs e Diagnóstico}

Em tarefas automáticas, os erros nem sempre aparecem na tela. Por isso, os logs são uma parte essencial do trabalho com \texttt{systemd}.

\begin{cmdbox}[Diagnóstico de Problemas]
\begin{tabularx}{\linewidth}{@{}C >{\raggedright\arraybackslash}X >{\raggedright\arraybackslash}X@{}}
\toprule
\textbf{Pergunta} & \textbf{Comando Útil} & \textbf{O que Observar} \\
\midrule
O serviço executou? & \texttt{systemctl status captura.service} & Estado, código de saída e mensagens \\
O timer está ativo? & \texttt{systemctl list-timers} & Próxima execução e última execução \\
Houve erro no script? & \texttt{journalctl -u captura.service -n 30} & Mensagens do script e falhas \\
As imagens foram geradas? & \texttt{ls -lh imagens} & Arquivos com data e hora \\
\bottomrule
\end{tabularx}
\end{cmdbox}

Problemas comuns incluem caminho incorreto, falta de permissão de execução, usuário errado no serviço, câmera indisponível ou ausência do comando \texttt{fswebcam}.

\section{Coleta Contínua com Apenas um Serviço}

O modelo com \texttt{service + timer} é adequado quando uma tarefa pode ser executada uma vez e terminar. Porém, para coletas muito frequentes, pode ser interessante usar um script contínuo, que fica rodando em loop.

Nesse segundo modelo, não usamos timer. O próprio script controla o intervalo entre capturas usando \texttt{sleep}.

\subsection{Comparação entre os Modelos}

\begin{cmdbox}[Timer versus Serviço Contínuo]
\begin{tabularx}{\linewidth}{@{}C >{\raggedright\arraybackslash}X >{\raggedright\arraybackslash}X@{}}
\toprule
\textbf{Modelo} & \textbf{Como Funciona} & \textbf{Quando Usar} \\
\midrule
\texttt{service + timer} & O timer chama o serviço repetidamente & Intervalos moderados, como minutos \\
\texttt{service contínuo} & O serviço inicia um script que permanece em execução & Intervalos curtos, como segundos \\
\bottomrule
\end{tabularx}
\end{cmdbox}

\subsection{Script de Captura Contínua}

O script abaixo captura imagens continuamente e aguarda 15 segundos entre uma captura e outra.

\begin{cmdbox}[Script \texttt{captura\_continua.sh}]
\begin{tabularx}{\linewidth}{@{}C >{\raggedright\arraybackslash}X@{}}
\toprule
\textbf{Linha} & \textbf{Código} \\
\midrule
1 & \texttt{\#!/bin/bash} \\
2 & \texttt{} \\
3 & \texttt{mkdir -p /home/g0/imagens\_continuas} \\
4 & \texttt{} \\
5 & \texttt{while true; do} \\
6 & \texttt{\ \ \ ts=\$(date +\%Y\%m\%d\_\%H\%M\%S)} \\
7 & \texttt{\ \ \ arquivo="/home/g0/imagens\_continuas/foto\_\$\{ts\}.jpg"} \\
8 & \texttt{} \\
9 & \texttt{\ \ \ echo "Capturando imagem em \$arquivo"} \\
10 & \texttt{\ \ \ /usr/bin/fswebcam -r 1280x720 "\$arquivo"} \\
11 & \texttt{} \\
12 & \texttt{\ \ \ sleep 15} \\
13 & \texttt{done} \\
\bottomrule
\end{tabularx}
\end{cmdbox}

Como esse script não termina sozinho, o teste manual precisa ser interrompido com \texttt{Ctrl+C}.

\subsection{Serviço Contínuo}

Para um script contínuo, usamos um serviço do tipo \texttt{simple}. O processo iniciado pelo serviço permanece ativo enquanto o script estiver rodando.

\begin{cmdbox}[Arquivo \texttt{/etc/systemd/system/captura-continua.service}]
\begin{tabularx}{\linewidth}{@{}C >{\raggedright\arraybackslash}X@{}}
\toprule
\textbf{Linha} & \textbf{Código} \\
\midrule
1 & \texttt{[Unit]} \\
2 & \texttt{Description=Captura continua de imagens com fswebcam} \\
3 & \texttt{After=network.target} \\
4 & \texttt{} \\
5 & \texttt{[Service]} \\
6 & \texttt{Type=simple} \\
7 & \texttt{User=g0} \\
8 & \texttt{WorkingDirectory=/home/g0} \\
9 & \texttt{ExecStart=/home/g0/captura\_continua.sh} \\
10 & \texttt{Restart=always} \\
11 & \texttt{} \\
12 & \texttt{[Install]} \\
13 & \texttt{WantedBy=multi-user.target} \\
\bottomrule
\end{tabularx}
\end{cmdbox}

\begin{cmdbox}[Controle do Serviço Contínuo]
\begin{tabularx}{\linewidth}{@{}C >{\raggedright\arraybackslash}X >{\raggedright\arraybackslash}X@{}}
\toprule
\textbf{Comando} & \textbf{Função} & \textbf{Exemplo de Uso} \\
\midrule
\texttt{start} & Iniciar a coleta contínua & \texttt{sudo systemctl start captura-continua.service} \\
\texttt{status} & Verificar se está rodando & \texttt{systemctl status captura-continua.service} \\
\texttt{journalctl} & Ver mensagens da coleta & \texttt{journalctl -u captura-continua.service -n 30} \\
\texttt{stop} & Parar a coleta contínua & \texttt{sudo systemctl stop captura-continua.service} \\
\bottomrule
\end{tabularx}
\end{cmdbox}

\section{Cuidados Importantes}

Ao trabalhar com \texttt{systemd}, pequenos detalhes podem impedir a automação de funcionar. A tabela abaixo resume cuidados úteis durante a prática.

\begin{cmdbox}[Checklist de Automação]
\begin{tabularx}{\linewidth}{@{}C >{\raggedright\arraybackslash}X >{\raggedright\arraybackslash}X@{}}
\toprule
\textbf{Cuidado} & \textbf{Motivo} & \textbf{Como Conferir} \\
\midrule
Testar manualmente & Isola problemas do script & \texttt{./captura.sh} \\
Usar caminho absoluto & O serviço não depende do terminal & \texttt{/home/g0/captura.sh} \\
Dar permissão de execução & O script precisa poder rodar & \texttt{chmod +x captura.sh} \\
Recarregar o systemd & Arquivos novos precisam ser lidos & \texttt{sudo systemctl daemon-reload} \\
Verificar logs & Erros podem não aparecer na tela & \texttt{journalctl -u captura.service} \\
Parar serviços contínuos & Evita capturas indefinidas & \texttt{systemctl stop ...} \\
\bottomrule
\end{tabularx}
\end{cmdbox}

Antes de pedir ajuda, é recomendável verificar três pontos: se o script roda manualmente, se o serviço apresenta erro no \texttt{status} e se o \texttt{journalctl} mostra alguma mensagem relacionada a caminho, permissão ou câmera.

\section{Conclusão}

Neste módulo, vimos como o \texttt{systemd} pode ser usado para automatizar rotinas de captura de imagens em um sistema Linux. A partir de um script Bash, criamos um serviço para executar uma tarefa e um timer para repetir essa tarefa em intervalos definidos.

Também vimos uma alternativa para coletas rápidas: um serviço contínuo que inicia um script com loop e \texttt{sleep}. A diferença entre os modelos é fundamental: no primeiro, o timer controla o intervalo; no segundo, o intervalo está dentro do próprio script.

Em metrologia por imagem, essa automação ajuda a reduzir erros manuais, organizar dados e melhorar a repetibilidade da coleta. Com serviços, timers e logs, o experimento deixa de depender apenas da ação direta do operador e passa a ser controlado por uma rotina configurada no sistema.

\end{document}
